\centerline{\bf \Huge Графчики}
\centerline{\bf \Huge }

\begin{flushright}
        \scriptsize{Понимать не обязательно. Достаточно знать.\\Антиспиральщик. Гуррен Лаганн.}
\end{flushright}

\problem{1}
В особой больнице есть главврач и много пациентов. В течение недели каждый пациент один раз в день кусал кого-нибудь (возможно и себя). В конце недели оказалось, что у всех пациентов по 2 укуса, а у главврача - 100 укусов. Сколько в больнице пациентов?

\problem{2}
Каждый мальчик знаком с 7 мальчиками и 10 девочками, а каждая девочка - с 8 девочками и 9 мальчиками. Кого больше - мальчиков или девочек?

\problem{3}
В углах доски $3 \times 3$ стоят кони: по 2 коня черного и белого цветов, причем кони одного цвета стоят в противоположных углах. За один ход разрешается выбрать любого коня и сделать им ход в свободную клетку. Можно ли переставить коней так, чтобы по прежнему все кони стояли в углах, но кони одного цвета стояли в углах при одной стороне?

\problem{4}
Дима выписал на доску первые 20 простых чисел: $2,3,5,7,11, \ldots$. Серёжа хочет соединить эти числа рёбрами так, чтобы из каждого числа выходило столько рёбер, какова величина этого числа. Два числа могут быть соединены несколькими рёбрами. Удастся ли Серёже осуществить желаемое?

\problem{5}
В королевстве живут графы, герцоги и маркизы. Однажды каждый граф сразился на дуэли с тремя герцогами и несколькими маркизами. Каждый герцог сразился на дуэли с двумя графами и шестью маркизами. Каждый маркиз сразился на дуэли с тремя герцогами и двумя графами. Известно, что все графы сразились с равным числом маркизов. Со сколькими маркизами сразился каждый граф?

\problem{6}
Футбольный мяч шьётся из 32 кусочков кожи: белых шестиугольников и чёрных пятиугольников. Каждый чёрный кусочек граничит только с белыми кусочками, каждый белый кусочек граничит с тремя чёрными и тремя белыми. Сколько чёрных кусочков нужно для изготовления мяча?